%!TEX root = ../main.tex
\section{Evaluation and Cross-Cutting Synthesis}
\label{sec:evaluation}

The literature contains many positive within-paper results but few fair
cross-paper comparisons.  Robots, tasks, data pools, policy families, model
initializations, training steps, extra pretraining data, and success definitions
usually differ.  A number reported in one setting establishes that an operation
can work under that setting; it does not rank the operation against a number from
another paper.  This section specifies the evidence needed for comparison and
then synthesizes conclusions that recur despite protocol differences.

\subsection{Three Levels of Outcome}

\textbf{Intrinsic outcomes} describe the prepared data: invalid-file rate,
label-noise estimate, duplicate rate, smoothness, entropy, coverage, balance, or
retained fraction.  They are cheap, interpretable diagnostics.  They do not show
that the measured property is causally useful to a policy.

\textbf{Proxy outcomes} include foundation-model scores, information estimates,
influence, predicted executability, inverse-kinematics feasibility, simulator
replay, or a learned success classifier.  They are closer to a target, but their
calibration set, model bias, and operating point must be reported.  The same proxy
should not define both the filter and the test of filtering.

\textbf{Downstream outcomes} retrain a policy and measure success, generalization,
sample efficiency, recovery, and safety.  This is the strongest evidence for
training utility and the most expensive.  Simulation affords repetitions and
factor control; real execution tests physical validity but often has few trials.
The two should be reported separately, with means, uncertainty, and trial counts.
Evidence is stronger when an effect persists across learners, data budgets,
embodiments, and held-out task factors.

\begin{table*}[t]
\centering
\caption{Minimum controls for evaluating a data-preparation intervention.}
\label{tab:evaluation-protocol}
\small
\setlength{\tabcolsep}{5pt}
\renewcommand{\arraystretch}{1.08}
\begin{tabularx}{\textwidth}{@{}p{2.6cm}X X@{}}
\toprule
Component & Required control & Failure prevented \\
\midrule
Data accounting & Raw and retained units, duration, source mix, lineage, leakage audit & Easier subsets presented as intrinsically better \\
Training budget & Same learner, initialization, updates, samples, augmentation, and tuning budget & Gains caused by extra optimization \\
Diagnostic and proxy & Defect labels, calibration split, uncertainty, threshold, confounders & Proxy scores treated as policy utility \\
Downstream utility & Multiple seeds, intervals, familiar, compositional, and shifted tasks & Rankings based on noisy point estimates \\
Closed-loop evidence & Separate simulation and robot trials, with logged failures and interventions & Evidence with different physical validity merged \\
Cost and risk & Human time, model and training compute, robot hours, discarded data, safety events & Preparation cost externalized \\
Robustness & Contamination, episode length, learner, target set, and budget sensitivity & Conclusions tied to one setting \\
\bottomrule
\end{tabularx}
\end{table*}


\subsection{Controlling the Intervention}

Table~\ref{tab:evaluation-protocol} gives a minimum protocol.  Its central design
is a factorial comparison among: unprepared full data; a random subset matched to
the retained volume; the proposed prepared subset; and, where relevant, an oracle
using ground-truth defect or utility labels.  A volume-matched random control
separates selection from regularization or fewer updates.  The full-data control
tests whether information was unnecessarily removed.  The oracle quantifies how
much headroom is due to the metric versus the underlying intervention.

Compute must be controlled in two complementary regimes.  Under
\emph{fixed-update} evaluation, every policy receives the same number of gradient
steps and batch samples, measuring value per training token/transition.  Under
\emph{fixed-pass} evaluation, each sees the same number of passes through its own
data, measuring an end-to-end recipe.  Reporting both reveals whether a smaller
curated set helps because its units are more valuable or because it is revisited
more often.  Hyperparameter tuning budgets must also match: a method-specific
threshold tuned on the test tasks is extra supervision.

Data leakage is unusually subtle in embodied settings.  Near-identical scene
configurations can occur in different episodes; generated variants can cross a
split; a VLM labeler may have seen benchmark images; and target demonstrations
used for influence can overlap evaluation objects.  Splits should occur before
augmentation and at the highest causal grouping available (task instance, object,
scene, collection session, and operator), with lineage-based duplicate checks.

\subsection{Cost, Openness, and Reproducibility}

``Using 50\% of the data'' is not a cost result if choosing that half required
training many policies, running target-robot rollouts, or manually labeling every
episode.  Report preparation cost in human minutes, foundation-model calls,
accelerator hours, simulator rollouts, real-robot hours, storage/I/O, and safety
interventions.  Then report total cost to reach a target deployment utility.
Expensive learner-aware selection can still be worthwhile when amortized across
many long pretraining runs; a cheap rule can be preferable for a small adaptation
run.

Reproducibility requires more than releasing final data.  The operator code,
model/checkpoint used for scoring, prompt and threshold, source version, exclusion
log, learned weights, random seeds, and lineage should be available.  OpenVLA and
Octo make policy baselines more accessible
~\cite{kim2024openvla_2406_09246,team2024octo_2405_12213}, but reproducible
preparation additionally needs a raw or defect-injected pool on which alternative
operators can be rerun.  Closed industrial reports can reveal scale and recipes,
yet claims depending on unavailable data or internal rollouts remain weaker
evidence for causal comparison.

\subsection{What Transfers from Language and Multimodal Data Engineering}

The most useful lesson from language-model data work is methodological: treat
the data pipeline itself as an optimizable, versioned system.  RefinedWeb and
FineWeb expose staged filtering and deduplication; DataComp-LM fixes a 240T-token
candidate pool, model/compute tracks, and 53 evaluations so that preparation can
be compared under a controlled training intervention
~\cite{penedo2023refinedweb_2306_01116,penedo2024fineweb_2406_17557,li2024datacomp_2406_11794}.
DoReMi and DSIR learn mixture or resampling decisions instead of treating source
ratios as invisible constants~\cite{xie2023doremi_2305_10429,xie2023data_2302_03169}.
LESS estimates optimizer-aware influence for targeted instruction tuning, while
DEITA separates quality, complexity, and diversity rather than selecting on one
score~\cite{xia2024less_2402_04333,liu2023what_2312_15685}.  JEST further shows
that the useful object can be a jointly selected batch, not an independently
ranked example~\cite{evans2024data_2406_17711}.  Table~\ref{tab:llm-lessons}
maps these mechanisms to embodied preparation and, equally importantly, records
what fails to transfer.

\begin{table*}[t]
\centering
\caption{Lessons from language and multimodal data engineering and the additional control required for embodied data.}
\label{tab:llm-lessons}
\scriptsize
\setlength{\tabcolsep}{4pt}
\begin{tabularx}{\textwidth}{@{}p{3.1cm}p{4.1cm}p{4.6cm}X@{}}
\toprule
Mechanism & Transferable lesson & Embodied adaptation & Additional control \\
\midrule
Fixed-pool competition~\cite{li2024datacomp_2406_11794,gadre2023datacomp_2304_14108}
& Hold the candidate pool, learner, compute, and evaluation fixed.
& Submit unit membership, weights, roles, lineage, and preparation cost.
& Separate simulation from blinded real-robot evaluation. \\
Auditable filtering~\cite{penedo2023refinedweb_2306_01116,penedo2024fineweb_2406_17557}
& Version each gate and deduplicate before splitting.
& Check technical validity, cross-modal alignment, and split leakage.
& Preserve temporal order and contact differences between visually similar episodes. \\
Mixture learning~\cite{xie2023doremi_2305_10429,xie2023data_2302_03169}
& Learn source weights against a declared target.
& Use hierarchical weights and protect rare but important groups.
& Test transfer from proxy learner or simulation to the target robot. \\
Targeted selection~\cite{xia2024less_2402_04333,liu2023what_2312_15685}
& Value depends on target, difficulty, and diversity.
& Select complementary units for a declared learner and task.
& Recalibrate after the policy changes the observed state distribution. \\
Joint batch selection~\cite{evans2024data_2406_17711}
& Example value depends on the surrounding batch.
& Select coherent skill or failure--recovery groups.
& Respect episode hierarchy and causal order. \\
Self-judged iteration~\cite{yuan2024self_2401_10020,shumailov2024collapse}
& Generation and judging can form a loop but may lose distribution tails.
& Retain audited real data while iterating on generated supervision.
& Use an independent verifier, mixture bounds, and rollback. \\
\bottomrule
\end{tabularx}
\end{table*}


\textbf{Pretraining suggests controlled pipeline competition.}
The DataComp-LM design is more important for EI than any winning text filter.
It answers a causal question: under a fixed raw pool, learner, token budget, and
evaluation suite, which preparation pipeline produces the best model?  Robot
papers often instead compare a curated set and a baseline with different episode
counts, updates, augmentations, or policy architectures.  A DataComp-Robot design
must preserve the fixed-intervention logic while adding action semantics,
episode/segment group splits, simulation and real tiers, preparation cost, and
safety.  Unlike text tokens, robot transitions are not freely exchangeable, and
their acquisition price depends on resets, supervision, embodiment, and risk.

\textbf{Mixture optimization suggests hierarchical, stage-specific routing.}
Language corpora have domains; embodied corpora have nested source, task,
environment, skill, embodiment, episode, segment, and transition groups.  A
single learned source weight is therefore only the top of the routing hierarchy.
A useful transfer of DoReMi/DSIR is to train a cheap proxy learner and estimate
which groups reduce target regret, then test whether weights transfer to a larger
VLA.  The non-transferable assumption is stationarity.  Once a robot policy is
deployed, its failures change and its data distribution follows; a ratio learned
offline becomes a controller prior, not a permanent optimum.  Minimum floors are
also needed for rare safety/recovery groups even when their short-term average
loss contribution is small.

\textbf{Instruction tuning suggests target-conditional selection below the
episode.}  LESS and DEITA support two ideas: value should be conditioned on a
target learner/task, and correctness, difficulty, and diversity should not be
collapsed prematurely.  For EI, the analogous candidate is often not a whole
episode but a coherent skill segment, correction, or failure--recovery pair.
Gradient influence is an informative proxy, yet a pointwise approximation can
miss temporal credit, complementary approach/grasp segments, and the
distribution shift caused by the improved policy.  The proper evaluation is
whether a selected hierarchical set reduces fixed-budget policy regret across
held-out tasks and learners, not whether its proxy score is high.

\textbf{LLM post-training suggests iterative judges, but also their failure
mode.}  Self-Rewarding Language Models illustrate a loop in which a model
generates candidates, judges responses, constructs preference pairs, and updates
itself~\cite{yuan2024self_2401_10020}.  The embodied analogue is attractive:
generate counterfactual failures, score trajectories with a reward VLM, collect
new behavior, and repeat.  It is also more dangerous.  A VLM may accept a
physically impossible contact, and a policy can change the visual distribution
on which its judge was calibrated.  Work on recursively generated training data
shows the more general risk: indiscriminate replacement of natural data can
erase distribution tails across generations~\cite{shumailov2024collapse}.
We do not claim the language theorem directly predicts robot-policy collapse;
we infer a design constraint.  Self-evolving embodied pipelines should
accumulate audited real anchors, identify synthetic lineage, use independent
verifiers, enforce subgroup floors, and test each round for tail and safety
retention.

The mapping is therefore \emph{mechanism-level}, not heuristic-level.  Text
pipelines teach how to version operators, run fixed-compute ablations, optimize
mixtures, estimate target influence, and audit recursive generation.  They do
not supply action causality, timing, controller semantics, physical validity, or
safe exploration.  Those missing properties explain why autonomous embodied
data preparation requires the layered evidence and control architecture in
Section~\ref{sec:recipes-engines}, rather than a direct port of web-data filters
or an LLM-as-judge.

\subsection{A DataComp-Robot Evaluation Track}

We propose a benchmark with four assets.  First, a versioned raw pool contains
multiple embodiments and sources, including natural and controlled defects, failed
and recovery episodes, redundant coverage, and aligned test leakage traps.
Second, a preparation submission emits training units, weights, roles, lineage,
and cost logs rather than a single cleaned archive.  Third, fixed lightweight and
foundation-policy learners are trained at multiple update budgets.  Fourth,
evaluation separates intrinsic defect detection, held-out simulation, and a
smaller blinded real-robot audit across in-distribution, compositional, and shifted
conditions.  A Pareto leaderboard reports deployment utility, worst-group
performance, preparation/training cost, and safety---not one scalar rank.

The benchmark should include intervention stress tests: contamination level,
episode length, label noise, action delay, rare recovery frequency, source
imbalance, and embodiment mismatch.  A method that detects one injected artifact
but fails under natural mixed defects should not be called a general quality
metric.  Fixed raw snapshots make results repeatable; periodic hidden refreshes
measure overfitting and support a living survey.

\subsection{Eight Cross-Paper Principles}

\textbf{P1: quality is target-conditional.}  Formal data-quality analysis,
mixture learning, datamodels, and influence methods all condition value on a
learner or target, explicitly or implicitly
~\cite{belkhale2023data_2306_02437,hejna2024re_2408_14037,dass2025datamil_2505_09603,lee2026quality_2603_09056}.
A universal score is therefore best interpreted as a prior, later recalibrated by
target evidence.

\textbf{P2: correctness and coverage form a frontier.}  Strong filtering removes
harmful labels and also removes rare states.  The mutual-information hard-task
result, the curation-metric audit, and Re-Mix's abnormal-source failure show why a
single threshold or learned weight can move to the wrong point on that frontier
~\cite{hejna2025robot_2502_08623,bedi2026what_2606_10229,hejna2024re_2408_14037}.

\textbf{P3: the actionable unit is moving below the episode.}  Segment recovery,
progress-based selection, retry value, and transition influence extract useful
behavior from mixed-quality trajectories
~\cite{wu2024learning_2401_08957,zhang2025scizor_2505_22626,wu2026sieve_2607_06442,qin2026beyond_2606_24633,lee2026quality_2603_09056}.
This shift requires uncertainty-aware boundaries and hierarchical lineage.

\textbf{P4: cross-embodiment standardization must align consequences, not shapes.}
Common schemas enable scale, but spatial frame, control quantity, timing,
morphology, and controller dynamics determine whether a label means the same
transition.  Structural masks are necessary interfaces, not complete semantic
alignment.

\textbf{P5: generation moves the bottleneck to verification.}  Segment
recomposition and world-model synthesis can rapidly expand nominal data
~\cite{mandlekar2023mimicgen_2310_17596,li2026xiaomi_2607_11643}.
The limiting evidence is whether semantics, geometry, kinematics, contact, and
policy utility survive jointly.

\textbf{P6: failed data is valuable when structured.}  An unlabeled failed action
sequence is not equivalent to a marked retry, valid prefix, recovery, negative,
or safety boundary.  Conflicting empirical claims usually arise from these
different semantics, not from a universal benefit or harm of failure.

\textbf{P7: data and models must be co-versioned.}  In a feedback loop, the model
chooses data that trains its successor.  Reproducing or rolling back a result
requires the raw source, preparation policy, verifier, training role, model
version, and deployment evidence as one causal graph.

\textbf{P8: self-evolution is a constrained experiment, not data accumulation.}
RECAP, SOP, and LWD obtain relevance from current-policy experience, but each
also retains an offline/pretrained reference or limits the update path
~\cite{intelligence20250_2511_14759,pan2026sop_2601_03044,wang2026learning_2605_00416}.
Recursive-data results explain why this pattern is principled rather than merely
conservative~\cite{shumailov2024collapse}: each round needs a hypothesis about a
gap, a controlled promotion test, anchor retention, and a rollback condition.
More self-generated data is not evidence of positive evolution.

\subsection{Paradigm Evolution}

The field's evolution is better understood by increasing data-decision authority
than by year (Figure~\ref{fig:autonomy-ladder}).  Static rule pipelines first
made records technically usable.  Pretrained models then supplied scalable
semantic labels and verifiers.  Learner-aware selectors used loss, datamodels,
influence, and progress to estimate target-conditional value.  Failure-driven
post-training added policy-induced states and interventions.  Current data
engines repeatedly couple acquisition, verification, routing, and deployment;
the prospective preparation agent would also choose the experiment and know
when to abstain, stop, or roll back.  Each step raises target relevance and
decision scope, but also increases cost, endogeneity, and feedback risk.  A
stable architecture is consequently a cascade---rules for obvious defects,
models for scalable semantics, learner-aware estimates for disputed allocation,
and sparse simulation/real audits for calibration---inside a versioned
controller, not a wholesale replacement of simple checks by one autonomous
judge.
